%% ma: L12-B08-0051
%% lop: 12
%% chuong: 2
%% bai: 8
%% dang: 12.8.39
%% loai: DS
%% muc: [TH, TH, TH, VD]
%% dap_an: "ĐĐĐĐ"
%% nguon: "LaTeX giáo viên gửi 21/10/2026 (M05, Vectơ và hệ trục tọa độ trong không gian, 107 câu), Phần II Câu 3"
%% kien_thuc: "Bài 7 (tọa độ điểm), Bài 8 (trung điểm, khoảng cách hai điểm)"
%% trang_thai: chinh-thuc
%% ghi_chu: "Giáo viên đã duyệt 21/10/2026. Đề gốc ghi 'cách điểm xuất phát' (chưa nói điểm xuất phát), đã thay bằng 'cách O' và bổ sung đơn vị km cho hệ trục; bỏ tiêu đề '(Bài toán Radar không lưu)'"
\begin{ex}
Xét hệ trục tọa độ $Oxyz$ (đơn vị: km) gắn với đài kiểm soát không lưu tại $O$, mặt phẳng $(Oxy)$ là mặt đất ngang, trục $Ox$ hướng về phía Nam, trục $Oy$ hướng về phía Đông, trục $Oz$ hướng thẳng đứng lên trời. Khinh khí cầu thứ nhất cách $O$ $2{,}5$ km về phía Nam, $2$ km về phía Đông và ở độ cao $0{,}8$ km. Khinh khí cầu thứ hai cách $O$ $1{,}5$ km về phía Bắc, $3$ km về phía Tây và ở độ cao $0{,}6$ km.
\choiceTF
{\True Tọa độ khinh khí cầu thứ nhất là $A(2{,}5;2;0{,}8)$.}
{\True Tọa độ khinh khí cầu thứ hai là $B(-1{,}5;-3;0{,}6)$.}
{\True Tọa độ trung điểm của đoạn thẳng nối hai khinh khí cầu là $I(0{,}5;-0{,}5;0{,}7)$.}
{\True Khoảng cách theo đường chim bay giữa hai khinh khí cầu xấp xỉ $6{,}4$ km.}
\loigiai{a) Nam là chiều dương $Ox$, Đông là chiều dương $Oy$: $A(2{,}5;2;0{,}8)$. Đúng.

b) Bắc là chiều âm $Ox$, Tây là chiều âm $Oy$: $B(-1{,}5;-3;0{,}6)$. Đúng.

c) $I\left(\dfrac{2{,}5-1{,}5}2;\dfrac{2-3}2;\dfrac{0{,}8+0{,}6}2\right)=(0{,}5;-0{,}5;0{,}7)$. Đúng.

d) $\overrightarrow{AB}=(-4;-5;-0{,}2)$, $AB=\sqrt{16+25+0{,}04}\approx6{,}41$ (km). Đúng.}
\end{ex}
